% Multi-chapter book fixture for the semantic resolver.
%
% Exercises long-book conventions that the article-shaped fixtures don't:
%   - explicit \chapter{...} commands (with hyperref labels)
%   - chapter-relative figure / equation numbering ("Figure 1.1", "Eq. (2.1)")
%   - footnotes on multiple pages
%   - a single bibliography at the end
%   - cross-references that span chapters
%
% Compiles with `pdflatex multi_chapter_book.tex` (book class). LaTeXML
% can also process this without issue — `tools/run_semantic_benchmark.py`
% feeds the .tex through LaTeXML for ground-truth marker extraction.

\documentclass{book}
\usepackage{amsmath}
\usepackage{hyperref}

\title{A Small Book on Markers}
\author{pdf2md test suite}
\date{}

\begin{document}
\maketitle
\tableofcontents

\chapter{Foundations}\label{ch:foundations}

This chapter introduces the basic vocabulary that the rest of the book
builds on. It also has a footnote\footnote{Footnotes are numbered
per-chapter in the book class.}.

\section{Notation}\label{sec:notation}
We use bold-face for vectors and italics for scalars.

\section{An equation}\label{sec:eq}
The opening identity is
\begin{equation}\label{eq:opening}
e^{i\pi} + 1 = 0
\end{equation}
which appears again in \autoref{ch:applications}.

\chapter{Applications}\label{ch:applications}

This chapter applies the foundation from \autoref{ch:foundations} to
two examples\footnote{The first applied example was independently
shown by Smith~\cite{smith2020}.}.

\section{First example}\label{sec:first}
By \autoref{eq:opening} the simplest case reduces to the identity.

\begin{figure}[h]
\centering
\fbox{\rule{0pt}{0.8in}\rule{1.5in}{0pt}}
\caption{A boxed placeholder.}\label{fig:box}
\end{figure}

\section{Second example}\label{sec:second}
\autoref{fig:box} shows the placeholder. Compare with the discussion
in \autoref{sec:notation}.

\begin{equation}\label{eq:closing}
\sum_{n=1}^{N} n = \frac{N(N+1)}{2}
\end{equation}

\chapter{References}\label{ch:references}

This is intentionally short — the bibliography is below.

\begin{thebibliography}{9}
\bibitem{smith2020} J.~Smith, \emph{Examples Worth Looking At}, 2020.
\bibitem{doe2021} J.~Doe, \emph{More Examples}, 2021.
\end{thebibliography}

\end{document}
